%% Question bank chapter file - generated by bank_tools.py build (do not edit while the bank is still being built:
%% it is regenerated from the per-test data in Question Bank\data). Every question is one bankq block:
%%   \begin{bankq}[id=..., source={...}, paper=.., qno=.., type=..., subtopic={...}, answer={...}, answer_src=key|solved,
%%                 working={...}, difficulty=1-10, diff_parts={K S I C F}, diff_note={...},
%%                 group=<passage id or empty>, also_in={duplicate copies in other tests}]
%% difficulty = min(10, K+S+I+C+F-1): K concepts 1-3, S steps 1-4, I insight 0-2, C computation 0-2, F format 0-1
%%   ...question LaTeX, options, figures (\bfig / \bfigi -> img/)...
%%   \end{bankq}
%% Passages shared by several questions are bankpassage blocks; their questions carry group=<passage id>.
\documentclass[10pt]{article}
\usepackage[a4paper,margin=1.8cm]{geometry}
\usepackage[T1]{fontenc}
\usepackage{amsmath,amssymb,mathtools}
\usepackage{graphicx}
\usepackage[export]{adjustbox}
\usepackage{xcolor,needspace,keyval}
\usepackage{enumitem}
\usepackage{booktabs,array,tabularx}
\usepackage[hidelinks]{hyperref}
\setlength{\parindent}{0pt}
\setlength{\parskip}{3pt}
\makeatletter
\define@key{bq}{id}{\def\bq@id{#1}}\define@key{bq}{source}{\def\bq@src{#1}}\define@key{bq}{paper}{\def\bq@paper{#1}}
\define@key{bq}{qno}{\def\bq@qno{#1}}\define@key{bq}{type}{\def\bq@type{#1}}\define@key{bq}{subtopic}{\def\bq@sub{#1}}
\define@key{bq}{answer}{\def\bq@ans{#1}}\define@key{bq}{answer_src}{\def\bq@asrc{#1}}\define@key{bq}{working}{\def\bq@work{#1}}
\define@key{bq}{group}{\def\bq@group{#1}}\define@key{bq}{also_in}{\def\bq@also{#1}}
\define@key{bq}{difficulty}{\def\bq@diff{#1}}\define@key{bq}{diff_parts}{\def\bq@dparts{#1}}\define@key{bq}{diff_note}{\def\bq@dnote{#1}}
\newenvironment{bankq}[1][]{\def\bq@work{}\def\bq@also{}\def\bq@group{}\def\bq@diff{?}\def\bq@dparts{}\def\bq@dnote{}\setkeys{bq}{#1}%
  \par\needspace{4\baselineskip}\medskip\noindent{\small\textbf{\bq@id}\hfill\color{gray}[\bq@src, Paper \bq@paper, Q\bq@qno]
  \bq@type\ $\cdot$ \bq@sub\ $\cdot$ \textbf{D\bq@diff}}\par\nopagebreak\relax}
  {\par{\small\color{teal}Answer: \textbf{\bq@ans}\ifx\bq@work\@empty\else\ (solved: \bq@work)\fi
  \ifx\bq@also\@empty\else\par\color{gray}Also in: \bq@also\fi}\par}
\define@key{bp}{id}{\def\bp@id{#1}}\define@key{bp}{source}{\def\bp@src{#1}}\define@key{bp}{questions}{\def\bp@qs{#1}}
\newenvironment{bankpassage}[1][]{\setkeys{bp}{#1}\par\needspace{4\baselineskip}\bigskip\noindent
  {\small\textbf{Passage \bp@id}\hfill\color{gray}[\bp@src; questions \bp@qs]}\par\nopagebreak\relax\itshape}{\par}
\makeatother
\newcommand{\bfig}[2]{\par\smallskip{\centering\includegraphics[width=#1pt,max width=0.62\linewidth,max height=0.32\textheight]{img/#2}\par}\smallskip}
\newcommand{\bfigi}[2]{\includegraphics[width=#1pt,max width=\linewidth,max height=0.32\textheight,valign=c]{img/#2}}
\newcommand{\bmissing}[1]{\textit{[figure missing in the source conversion: \detokenize{#1}]}}
\newenvironment{qoptions}{\begin{itemize}[leftmargin=1.8em,itemsep=1pt,topsep=2pt,label={}]}{\end{itemize}}
\newcommand{\opt}[2]{\item[(#1)] #2}
\providecommand{\fwstart}{}\providecommand{\fwend}{}\providecommand{\qfullsplit}{}
\newcolumntype{L}[1]{>{\raggedright\arraybackslash}p{#1}}
\newcolumntype{C}[1]{>{\centering\arraybackslash}p{#1}}
\newcolumntype{R}[1]{>{\raggedleft\arraybackslash}p{#1}}
\providecommand{\sectionhdr}[1]{}\providecommand{\subsectionhdr}[1]{}


\begin{document}
{\Large\bfseries Circles --- Question Bank}\par\smallskip
{\small Mathematics chapter M16. Questions from all converted test papers, grouped by source series. Answers: printed key where it matches the paper, otherwise solved (answer\_src=solved, with working). Exact repeats across papers appear once, with also\_in listing the other copies.}\par\medskip\hrule

\section*{Score-I papers}
\begin{bankq}[id=SI12-P1-MATH-S3-Q06, source={Score-I - Test 12 (09 Mar 2022)}, paper=1, qno=6, type=numerical, subtopic={Tangents}, answer={1}, answer_src=key, difficulty=5, diff_parts={K2 S2 I1 C0 F1}, diff_note={concepts: tangents from a point on the director circle, tangent circles | steps: tangents are perpendicular; centre on OP with distance r sqrt 2 from P; solve for r | I: use the symmetric position of the centre on OP}]
If $C_1:x^2+y^2=(3+2\sqrt2)^2$ be a circle and $PA$ and $PB$ are pair of tangents on $C_1$ where $P$ is any point on the director circle of $C_1$, then the radius of smallest circle which touches $C_1$ externally and also the two tangents $PA$ and $PB$, is --
\end{bankq}
\begin{bankq}[id=SI12-P1-MATH-S3-Q07, source={Score-I - Test 12 (09 Mar 2022)}, paper=1, qno=7, type=numerical, subtopic={Family of circles}, answer={1}, answer_src=key, difficulty=3, diff_parts={K2 S1 I0 C0 F1}, diff_note={concepts: locus parametrisation, eccentricity of an ellipse | steps: parametrise P and Q to get x^2/b^2 + y^2/a^2 = 1; foci on the inner circle gives b^2 = 2a^2 and e}]
Origin $O$ is the centre of two concentric circles whose radii are $a$ \& $b$ respectively, $a<b$. A line $OPQ$ is drawn to cut the inner circle in $P$ \& the outer circle in $Q$. $PR$ is drawn parallel to the $y$-axis \& $QR$ is drawn parallel to the $x$-axis. The locus of $R$ is an ellipse touching the two circles and the focii of this ellipse lie on the inner circle. if eccentricity of that ellipse is $\lambda$, then find $\sqrt2\lambda$
\end{bankq}
\begin{bankpassage}[id=SI12-P2-MATH-S1-G7-8, source={Score-I - Test 12 (09 Mar 2022)}, questions={7-8}]
$P(a,b)$ is a point in 1$^{\text{st}}$ quadrant through which circles pass touching both the coordinate axes.\\
Column 1 : Coordinates of point $P$.\\
Column 2 : Sum of ordinates of centers of all such circles.\\
Column 3 : Distance between the centers of such circles.
\begin{center}\footnotesize
\begin{tabular}{|p{0.6cm}|p{1.4cm}|p{0.6cm}|p{1.4cm}|p{0.5cm}|p{1.8cm}|}
\hline
 & \centering Column 1 & & \centering Column 2 & & \centering Column 3\tabularnewline
\hline
(I) & $(2,3)$ & (i) & 8 & (P) & $4\sqrt3$\\
(II) & $(3,1)$ & (ii) & 12 & (Q) & 8\\
(III) & $(2,2)$ & (iii) & 10 & (R) & $8\sqrt2$\\
(IV) & $(4,2)$ & (iv) & 6 & (S) & $4\sqrt6$\\
\hline
\end{tabular}
\end{center}
\end{bankpassage}
\begin{bankq}[id=SI12-P2-MATH-S1-Q07, source={Score-I - Test 12 (09 Mar 2022)}, paper=2, qno=7, type=single, subtopic={Family of circles}, answer={C}, answer_src=key, difficulty=3, diff_parts={K2 S2 I0 C0 F0}, diff_note={concepts: circles, tangents and family of circles | steps: circle for item I; circle for item II; circle for item III; circle for item IV and match}, group=SI12-P2-MATH-S1-G7-8]
Identify the correct combination ?
\begin{qoptions}
\opt{A}{(I) (iii) (P)}
\opt{B}{(II) (i) (R)}
\opt{C}{(III) (i) (Q)}
\opt{D}{(IV) (ii) (S)}
\end{qoptions}
\end{bankq}
\begin{bankq}[id=SI12-P2-MATH-S1-Q08, source={Score-I - Test 12 (09 Mar 2022)}, paper=2, qno=8, type=single, subtopic={Family of circles}, answer={A}, answer_src=key, difficulty=3, diff_parts={K2 S2 I0 C0 F0}, diff_note={concepts: circles, tangents and family of circles | steps: circle for item I; circle for item II; circle for item III; circle for item IV and match}, group=SI12-P2-MATH-S1-G7-8]
Identify the correct combination ?
\begin{qoptions}
\opt{A}{(I) (iii) (S)}
\opt{B}{(II) (i) (Q)}
\opt{C}{(III) (i) (R)}
\opt{D}{(IV) (ii) (P)}
\end{qoptions}
\end{bankq}

\section*{TOAS papers}
\begin{bankq}[id=TOAS06-P1-MATH-S1-Q03, source={TOAS - Test 06 (17 Apr 2022)}, paper=1, qno=3, type=single, subtopic={Tangents}, answer={C}, answer_src=key, difficulty=2, diff_parts={K1 S2 I0 C0 F0}, diff_note={concepts: circumscribed circle and tangents | steps: vertices of the square; circle x^2 + y^2 = 2; coordinates of P; PA^2 + PC^2}]
A circle is drawn circumscribing the square with vertices A, B, C, D (in order) represented by $|x+y|+|x-y|=2$. If tangents at A and B to this circle meets at P, then $\mathrm{PA}^2+\mathrm{PC}^2$ is equal to
\begin{qoptions}
\opt{A}{4}
\opt{B}{8}
\opt{C}{12}
\opt{D}{20}
\end{qoptions}
\end{bankq}
\begin{bankq}[id=TOAS06-P1-MATH-S2-Q06, source={TOAS - Test 06 (17 Apr 2022)}, paper=1, qno=6, type=multi, subtopic={Family of circles}, answer={A,B}, answer_src=key, difficulty=3, diff_parts={K1 S2 I0 C0 F1}, diff_note={concepts: circle through given tangent and normal conditions | steps: centre on the normal at (1,1); centre on the second line; radius; options on tangent length, intercepts}]
A circle is tangent to line $3x-4y+1=0$ at $(1,1)$ and is normal to line $2x+y-5=0$. Choose correct option(s)?
\begin{qoptions}
\opt{A}{Length of tangent to circle from $(-1,1)$ is 4 units}
\opt{B}{Length of perpendicular from centre of circle to $x$-axis is 3 units.}
\opt{C}{$x$-intercept of director circle of given circle is $\sqrt{82}$}
\opt{D}{$x$-intercept of director circle of given circle is $\sqrt{41}$}
\end{qoptions}
\end{bankq}
\begin{bankq}[id=TOAS06-P1-MATH-S3-Q05, source={TOAS - Test 06 (17 Apr 2022)}, paper=1, qno=5, type=numerical, subtopic={Tangents}, answer={1.87 or 1.88}, answer_src=key, difficulty=6, diff_parts={K2 S2 I1 C1 F1}, diff_note={concepts: tangent circles, a common diameter | steps: centres of the given circles; unknown radius in terms of the position; two tangency conditions; solve | I: choose coordinates along the diameter}]
Given two circle $x^2+y^2+3\sqrt2(x+y)=0$ and $x^2+y^2+5\sqrt2(x+y)=0$, the radius of third circle which is tangent to the two given circles and to their common diametre is
\end{bankq}
\end{document}
